\section{Computing intersection multiplicities}
\label{section-computing-intersection-multiplicities}

\noindent
In this section we discuss some cases where the intersection multiplicities
can be computed by different means. Here is a first example.

\begin{lemma}
\label{lemma-intersection-multiplicity-CM}
Let $X$ be a nonsingular variety and $W, V \subset X$ closed
subvarieties which intersect properly. Let $Z$ be an irreducible component
of $V \cap W$ with generic point $\xi$. Assume that $\mathcal{O}_{W, \xi}$
and $\mathcal{O}_{V, \xi}$ are Cohen-Macaulay. Then
$$
e(X, V \cdot W, Z) =
\text{length}_{\mathcal{O}_{X, \xi}}(\mathcal{O}_{V \cap W, \xi})
$$
where $V \cap W$ is the scheme theoretic intersection.
In particular, if both $V$ and $W$ are Cohen-Macaulay, then
$V \cdot W = [V \cap W]_{\dim(V) + \dim(W) - \dim(X)}$.
\end{lemma}

\begin{proof}
Set $A = \mathcal{O}_{X, \xi}$, $B = \mathcal{O}_{V, \xi}$, and
$C = \mathcal{O}_{W, \xi}$. By Auslander-Buchsbaum
(Algebra, Proposition \ref{algebra-proposition-Auslander-Buchsbaum})
we can find a finite free resolution $F_\bullet \to B$ of
length
$$
\text{depth}(A) - \text{depth}(B) =
\dim(A) - \dim(B) = \dim(C)
$$
First equality as $A$ and $B$ are Cohen-Macaulay and the second
as $V$ and $W$ intersect properly. Then $F_\bullet \otimes_A C$ is a
complex of finite free modules representing $B \otimes_A^\mathbf{L} C$
hence has cohomology modules with support in $\{\mathfrak m_A\}$.
By the Acyclicity lemma (Algebra, Lemma \ref{algebra-lemma-acyclic})
which applies as $C$ is Cohen-Macaulay
we conclude that $F_\bullet \otimes_A C$ has nonzero
cohomology only in degree $0$. This finishes the proof.
\end{proof}

\begin{lemma}
\label{lemma-one-ideal-ci}
Let $A$ be a Noetherian local ring. Let $I = (f_1, \ldots, f_r)$ be an ideal
generated by a regular sequence. Let $M$ be a finite $A$-module. Assume that
$\dim(\text{Supp}(M/IM)) = 0$. Then
$$
e_I(M, r) = \sum (-1)^i\text{length}_A(\text{Tor}_i^A(A/I, M))
$$
Here $e_I(M, r)$ is as in Remark \ref{remark-trivial-generalization}.
\end{lemma}

\begin{proof}
Since $f_1, \ldots, f_r$ is a regular sequence the Koszul complex
$K_\bullet(f_1, \ldots, f_r)$ is a resolution of $A/I$ over $A$, see
More on Algebra, Lemma
\ref{more-algebra-lemma-noetherian-finite-all-equivalent}.
Thus the right hand side is equal to
$$
\sum (-1)^i\text{length}_A H_i(K_\bullet(f_1, \ldots, f_r) \otimes_A M)
$$
Now the result follows immediately from
Theorem \ref{theorem-multiplicity-with-koszul} if $I$ is an ideal
of definition. In general, we replace $A$ by $\overline{A} = A/\text{Ann}(M)$
and $f_1, \ldots, f_r$ by $\overline{f}_1, \ldots, \overline{f}_r$
which is allowed because
$$
K_\bullet(f_1, \ldots, f_r) \otimes_A M =
K_\bullet(\overline{f}_1, \ldots, \overline{f}_r) \otimes_{\overline{A}} M
$$
Since $e_I(M, r) = e_{\overline{I}}(M, r)$ where
$\overline{I} = (\overline{f}_1, \ldots, \overline{f}_r) \subset \overline{A}$
is an ideal of definition the result follows from
Theorem \ref{theorem-multiplicity-with-koszul} in this case as well.
\end{proof}

\begin{lemma}
\label{lemma-multiplicity-with-lci}
Let $X$ be a nonsingular variety. Let $W,V \subset X$ be
closed subvarieties which intersect properly. Let $Z$ be an irreducible
component of $V \cap W$ with generic point $\xi$.
Suppose the ideal of $V$ in $\mathcal{O}_{X, \xi}$ is cut out by
a regular sequence $f_1, \ldots, f_c \in \mathcal{O}_{X, \xi}$.
Then $e(X, V\cdot W, Z)$ is equal to $c!$ times the leading coefficient in
the Hilbert polynomial
$$
t \mapsto \text{length}_{\mathcal{O}_{X, \xi}}
\mathcal{O}_{W, \xi}/(f_1, \ldots, f_c)^t,\quad t \gg 0.
$$
In particular, this coefficient is $> 0$.
\end{lemma}

\begin{proof}
The equality
$$
e(X, V\cdot W, Z) = e_{(f_1, \ldots, f_c)}(\mathcal{O}_{W, \xi}, c)
$$
follows from the more general Lemma \ref{lemma-one-ideal-ci}.
To see that $e_{(f_1, \ldots, f_c)}(\mathcal{O}_{W, \xi}, c)$ is
$> 0$ or equivalently that $e_{(f_1, \ldots, f_c)}(\mathcal{O}_{W, \xi}, c)$
is the leading coefficient of the Hilbert polynomial
it suffices to show that the
dimension of $\mathcal{O}_{W, \xi}$ is $c$, because the degree of the
Hilbert polynomial is equal to the dimension by
Algebra, Proposition \ref{algebra-proposition-dimension}.
Say $\dim(V) = r$, $\dim(W) = s$, and $\dim(X) = n$. Then
$\dim(Z) = r + s - n$ as the intersection is proper. Thus
the transcendence degree of $\kappa(\xi)$ over $\mathbf{C}$ is
$r + s - n$, see Algebra, Lemma
\ref{algebra-lemma-dimension-prime-polynomial-ring}.
We have $r + c = n$ because $V$ is cut out by a regular sequence
in a neighbourhood of $\xi$, see
Divisors, Lemma \ref{divisors-lemma-Noetherian-scheme-regular-ideal}
and then Lemma \ref{lemma-pullback-by-regular-immersion}
applies (for example). Thus
$$
\dim(\mathcal{O}_{W, \xi}) = s - (r + s - n) = s - ((n - c) + s - n) = c
$$
the first equality by Algebra, Lemma
\ref{algebra-lemma-dimension-at-a-point-finite-type-field}.
\end{proof}

\begin{lemma}
\label{lemma-multiplicity-with-effective-Cartier-divisor}
In Lemma \ref{lemma-multiplicity-with-lci} assume that  $c = 1$, i.e., $V$
is an effective Cartier divisor. Then
$$
e(X, V \cdot W, Z) =
\text{length}_{\mathcal{O}_{X, \xi}}
(\mathcal{O}_{W, \xi}/f_1\mathcal{O}_{W, \xi}).
$$
\end{lemma}

\begin{proof}
In this case the image of $f_1$ in $\mathcal{O}_{W, \xi}$ is nonzero by
properness of intersection, hence a nonzerodivisor divisor. Moreover,
$\mathcal{O}_{W, \xi}$ is a Noetherian local domain of dimension $1$.
Thus
$$
\text{length}_{\mathcal{O}_{X, \xi}}
(\mathcal{O}_{W, \xi}/f_1^t\mathcal{O}_{W, \xi}) =
t \text{length}_{\mathcal{O}_{X, \xi}}
(\mathcal{O}_{W, \xi}/f_1\mathcal{O}_{W, \xi})
$$
for all $t \geq 1$, see Algebra, Lemma \ref{algebra-lemma-ord-additive}.
This proves the lemma.
\end{proof}

\begin{lemma}
\label{lemma-multiplicity-lci-CM}
In Lemma \ref{lemma-multiplicity-with-lci} assume that
the local ring $\mathcal{O}_{W, \xi}$ is Cohen-Macaulay. Then we
have
$$
e(X, V \cdot W, Z) =
\text{length}_{\mathcal{O}_{X, \xi}} (\mathcal{O}_{W, \xi}/
f_1\mathcal{O}_{W, \xi} + \ldots + f_c\mathcal{O}_{W, \xi}).
$$
\end{lemma}

\begin{proof}
This follows immediately from Lemma \ref{lemma-intersection-multiplicity-CM}.
Alternatively, we can deduce it from Lemma \ref{lemma-multiplicity-with-lci}.
Namely, by Algebra, Lemma \ref{algebra-lemma-reformulate-CM}
we see that $f_1, \ldots, f_c$ is a regular sequence in
$\mathcal{O}_{W, \xi}$. Then
Algebra, Lemma \ref{algebra-lemma-regular-quasi-regular} shows that
$f_1, \ldots, f_c$ is a quasi-regular sequence.
This easily implies the length of
$\mathcal{O}_{W, \xi}/(f_1, \ldots, f_c)^t$ is
$$
{c + t \choose c}
\text{length}_{\mathcal{O}_{X, \xi}} (\mathcal{O}_{W, \xi}/
f_1\mathcal{O}_{W, \xi} + \ldots + f_c\mathcal{O}_{W, \xi}).
$$
Looking at the leading coefficient we conclude.
\end{proof}


